\label{sec:conclusion}

This paper has developed a systematic, quantitative methodology for satisfying
all three Gingles preconditions using the \texttt{redist} algorithmic
redistricting system.  The methodology addresses each prong with a specific
quantitative measure grounded in established statistical methods:

\begin{itemize}
  \item \textbf{Prong 1} is quantified by VRASection's geographic alignment
    score, with ensemble-based confidence bounds.  The 0.5 alignment threshold
    corresponds to the empirically observed lower bound of judicially accepted
    majority-minority districts.

  \item \textbf{Prong 2} is quantified by WLS ecological regression with HC3
    standard errors, applied to contested primary elections to isolate racial
    from partisan cohesion.

  \item \textbf{Prong 3} is quantified by the post-\textit{Callais}
    WLS+HC3+Holm regression with a partisan control variable, supplemented by
    LOO robustness checks and bootstrap confidence intervals.
\end{itemize}

The worked example using Alabama AL-07 demonstrates that the methodology
produces clear, defensible results across all three prongs.  The racial
bloc-voting coefficient drops from 0.71 to 0.43 when the partisan control
is introduced---a substantial reduction that confirms the importance of the
post-\textit{Callais} disentanglement requirement---but remains large,
statistically significant, and stable, satisfying Prong 3 under the more
demanding post-\textit{Callais} standard.

Several limitations should be noted.  First, the methodology relies on
precinct-level ecological regression, which is subject to ecological inference
bias~\citep{king1997solution}.  In districts with limited racial heterogeneity
across precincts, the regression estimates will be imprecise.  Second, the
alignment score depends on the VRASection algorithm's identification of the
minority community of interest, which is itself a modeling choice.  Third, the
Holm correction controls family-wise error rate but may be conservative in
elections with high correlations; the Benjamini-Hochberg procedure may be
more appropriate in such settings.

The expert witness checklist in Section~7 provides a standardized workflow
that can be applied consistently across Section 2 cases, reducing the scope
of expert disagreement to the substantive questions---whether the racial
bloc-voting coefficient is large enough, whether the alignment score threshold
is appropriate---rather than to methodological choices that should be
resolved by reference to established statistical practice.

For courts and practitioners, the key contribution is the \textit{Callais}
implementation.  Prior to \textit{Callais}, Prong 3 analysis was commonly
reported without a partisan control variable, producing inflated racial
bloc-voting coefficients that overstated the racial component.
Post-\textit{Callais}, the partisan control is required.  The methodology
developed here provides a specific, reproducible implementation of that
requirement using standard econometric tools available in the
\texttt{redist-analysis} crate.
